% Volume X: Normative Conclusions
% Part XVIII. Freedom Under Causal Opacity
% Chapter 110: The Freedom Created by Law
% Spine: proof-spines/ch110-spine.md

\chapter{The Freedom Created by Law}
\label{ch:ch110}

The final chapter returns to the apparent absurdity with which the book began: why should an innocent person have to endure formal proof, hearings, and delay? Because no institution receives innocence or guilt directly from the world. It receives records, testimony, measurements, classifications, and arguments.

The inconvenience of law creates freedom when it prevents those representations from becoming coercive facts too quickly. Legal procedure is therefore not merely an obstacle placed between knowledge and action. It is part of the machinery by which knowledge capable of authorizing action is produced.

The book closes, then, with its governing formulation:

\begin{quote}
\emph{Paranoia is powerful because causal reality exceeds every available account of it. Privacy prevents the account from becoming total exposure. Publicity prevents authority from becoming total secrecy. Evidence prevents suspicion from becoming fact. Procedure prevents fact from becoming force without passing through contestable form. Freedom persists in the intervals maintained among these operations.}
\end{quote}
